\[ %%%%%%%%%%%%%%%%%%%%%%%%%%%%%%%%%%%%%%%%%%%%%%%%%%%%%%%%%%%%%%%%%%%%%%%%%%%%%%%% \subsection{Baseline Models} \label{sec:baselines} %%%%%%%%%%%%%%%%%%%%%%%%%%%%%%%%%%%%%%%%%%%%%%%%%%%%%%%%%%%%%%%%%%%%%%%%%%%%%%%% To evaluate the effectiveness of PHYSGEN, we compare it with representative models from several major categories of scientific machine learning. The selected baselines cover sequence models, graph neural simulators, physics- informed approaches, neural operators, and hybrid differentiable physics methods. The comparison is designed to evaluate whether PHYSGEN provides advantages in long-horizon prediction, physical consistency, and dynamical stability. %%%%%%%%%%%%%%%%%%%%%%%%%%%%%%%%%%%%%%%%%%%%%%%%%%%%%%%%%%%%%%%%%%%%%%%%%%%%%%%% \subsubsection{Data-Driven Sequence Models} %%%%%%%%%%%%%%%%%%%%%%%%%%%%%%%%%%%%%%%%%%%%%%%%%%%%%%%%%%%%%%%%%%%%%%%%%%%%%%%% \paragraph{Multi-Layer Perceptron (MLP).} The MLP baseline represents a purely data-driven approximation of the transition operator: \begin{equation} \hat{X}_{t+1} = f_\theta(X_t). \end{equation} Although computationally efficient, MLP models do not explicitly represent physical interactions or constraints. %%%%%%%%%%%%%%%%%%%%%%%%%%%%%%%%%%%%%%%%%%%%%%%%%%%%%%%%%%%%%%%%%%%%%%%%%%%%%%%% \paragraph{Long Short-Term Memory Network (LSTM).} LSTM models temporal dependencies using recurrent hidden states: \begin{equation} h_{t+1} = LSTM(X_t,h_t). \end{equation} This baseline evaluates whether temporal memory alone can capture nonlinear system evolution. %%%%%%%%%%%%%%%%%%%%%%%%%%%%%%%%%%%%%%%%%%%%%%%%%%%%%%%%%%%%%%%%%%%%%%%%%%%%%%%% \subsubsection{Graph Neural Network Baselines} %%%%%%%%%%%%%%%%%%%%%%%%%%%%%%%%%%%%%%%%%%%%%%%%%%%%%%%%%%%%%%%%%%%%%%%%%%%%%%%% \paragraph{Graph Neural Network (GNN).} A standard message passing GNN is used as a graph-based baseline: \begin{equation} h_i^{t+1} = \phi ( h_i^t, \sum_{j\in N(i)} \psi(h_i^t,h_j^t) ). \end{equation} Unlike PHYSGEN, the standard GNN does not incorporate explicit physical constraints or stability regularization. %%%%%%%%%%%%%%%%%%%%%%%%%%%%%%%%%%%%%%%%%%%%%%%%%%%%%%%%%%%%%%%%%%%%%%%%%%%%%%%% \paragraph{Graph Neural Simulator (GNS).} The Graph Neural Simulator baseline represents a state-of-the-art approach for learning physical simulations from graph representations. The model predicts particle or mesh evolution: \begin{equation} X_{t+1} = GNS_\theta(G_t). \end{equation} This comparison evaluates the contribution of physics guidance beyond purely graph-based simulation. %%%%%%%%%%%%%%%%%%%%%%%%%%%%%%%%%%%%%%%%%%%%%%%%%%%%%%%%%%%%%%%%%%%%%%%%%%%%%%%% \subsubsection{Physics-Informed Learning Baselines} %%%%%%%%%%%%%%%%%%%%%%%%%%%%%%%%%%%%%%%%%%%%%%%%%%%%%%%%%%%%%%%%%%%%%%%%%%%%%%%% \paragraph{Physics-Informed Neural Networks (PINNs).} PINNs incorporate governing equations through physics residual minimization: \begin{equation} \mathcal{L}_{PINN} = \mathcal{L}_{data} + \lambda \mathcal{L}_{PDE}. \end{equation} PINNs provide strong physical consistency but may suffer from scalability limitations for large dynamical systems. %%%%%%%%%%%%%%%%%%%%%%%%%%%%%%%%%%%%%%%%%%%%%%%%%%%%%%%%%%%%%%%%%%%%%%%%%%%%%%%% \paragraph{Physics-Constrained Neural Networks.} A physics-constrained neural network baseline is implemented by adding physical penalties to a standard neural architecture: \begin{equation} \mathcal{L} = \mathcal{L}_{data} + \lambda_p \mathcal{L}_{constraint}. \end{equation} This baseline evaluates the importance of integrating physics directly into the graph evolution mechanism. %%%%%%%%%%%%%%%%%%%%%%%%%%%%%%%%%%%%%%%%%%%%%%%%%%%%%%%%%%%%%%%%%%%%%%%%%%%%%%%% \subsubsection{Neural Operator Baselines} %%%%%%%%%%%%%%%%%%%%%%%%%%%%%%%%%%%%%%%%%%%%%%%%%%%%%%%%%%%%%%%%%%%%%%%%%%%%%%%% \paragraph{Fourier Neural Operator (FNO).} FNO learns mappings between infinite-dimensional function spaces: \begin{equation} \mathcal{G} : a(x) \rightarrow u(x). \end{equation} It provides a strong baseline for PDE forecasting problems. %%%%%%%%%%%%%%%%%%%%%%%%%%%%%%%%%%%%%%%%%%%%%%%%%%%%%%%%%%%%%%%%%%%%%%%%%%%%%%%% \paragraph{Deep Operator Network (DeepONet).} DeepONet approximates nonlinear operators using branch and trunk networks: \begin{equation} u(x) = \sum_i b_i(a)t_i(x). \end{equation} This evaluates whether operator learning approaches can achieve comparable long-horizon performance. %%%%%%%%%%%%%%%%%%%%%%%%%%%%%%%%%%%%%%%%%%%%%%%%%%%%%%%%%%%%%%%%%%%%%%%%%%%%%%%% \subsubsection{Hybrid Differentiable Physics Baseline} %%%%%%%%%%%%%%%%%%%%%%%%%%%%%%%%%%%%%%%%%%%%%%%%%%%%%%%%%%%%%%%%%%%%%%%%%%%%%%%% \paragraph{Differentiable Physics Model.} A hybrid differentiable physics baseline combines numerical simulation with learned correction terms: \begin{equation} X_{t+1} = F_{physics}(X_t) + \Delta_\theta(X_t). \end{equation} This baseline evaluates the advantage of learning physics-guided graph evolution rather than correcting traditional simulators. %%%%%%%%%%%%%%%%%%%%%%%%%%%%%%%%%%%%%%%%%%%%%%%%%%%%%%%%%%%%%%%%%%%%%%%%%%%%%%%% \subsubsection{PHYSGEN Model} %%%%%%%%%%%%%%%%%%%%%%%%%%%%%%%%%%%%%%%%%%%%%%%%%%%%%%%%%%%%%%%%%%%%%%%%%%%%%%%% The proposed PHYSGEN model combines the advantages of the previous approaches: \begin{equation} X_{t+1} = \Phi_\theta ( G_t, P_t, C_t ). \end{equation} Unlike existing baselines, PHYSGEN integrates: \begin{itemize} \item adaptive graph representation; \item physics-guided message passing; \item latent physical state evolution; \item multi-level physical constraints; \item stability-regularized optimization. \end{itemize} %%%%%%%%%%%%%%%%%%%%%%%%%%%%%%%%%%%%%%%%%%%%%%%%%%%%%%%%%%%%%%%%%%%%%%%%%%%%%%%% \subsubsection{Baseline Comparison Summary} %%%%%%%%%%%%%%%%%%%%%%%%%%%%%%%%%%%%%%%%%%%%%%%%%%%%%%%%%%%%%%%%%%%%%%%%%%%%%%%% The evaluated methods are summarized in Table~\ref{tab:baseline_models}. \begin{table}[h] \centering \caption{Comparison of baseline models and PHYSGEN characteristics.} \label{tab:baseline_models} \begin{tabular}{lcccc} \hline \textbf{Model} & \textbf{Graph} & \textbf{Physics} & \textbf{Long Horizon} & \textbf{Stability} \\ \hline MLP & No & No & Low & No \\ LSTM & No & No & Medium & No \\ GNN & Yes & No & Medium & No \\ GNS & Yes & Partial & High & Partial \\ PINN & No & Yes & Medium & Partial \\ FNO & No & Partial & High & Partial \\ DeepONet & No & Partial & High & Partial \\ Differentiable Physics & Partial & Yes & High & Yes \\ PHYSGEN & Yes & Yes & High & Yes \\ \hline \end{tabular} \end{table} %%%%%%%%%%%%%%%%%%%%%%%%%%%%%%%%%%%%%%%%%%%%%%%%%%%%%%%%%%%%%%%%%%%%%%%%%%%%%%%% \subsubsection{Summary} %%%%%%%%%%%%%%%%%%%%%%%%%%%%%%%%%%%%%%%%%%%%%%%%%%%%%%%%%%%%%%%%%%%%%%%%%%%%%%%% The baseline selection provides a comprehensive comparison across major scientific machine learning paradigms. By evaluating PHYSGEN against sequence models, graph simulators, physics- informed approaches, and neural operators, the experiments aim to demonstrate the contribution of physics-guided graph evolution for stable long-horizon prediction of nonlinear dynamical systems. \] 
